  \section{Línea temporal de vida}

\noindent El objetivo de esta sección será describir a grandes rasgos la línea temporal que habita un microcuásar. Partiendo de los requisitos, su origen, los estados durmientes y activos, y su fin. En particular, se dejará la mayor parte del estado activo a la sección siguiente, <<El ciclo de estados espectrales>>. 

Debido al flujo de la literatura, esta sección será más genérica a sistemas binarios de rayos X, sistemas binarios con agujeros negros, y centralmente, LMXB (Binarios de rayos X de baja masa, \textit{Low Mass X-ray Binaries}). El estudio de estos cuerpos es una superclase al estudio de los microcuásarses, y por tanto no deja de ser relevante.

\begin{figure}[htbp]
\centering
\begin{tikzpicture}[
  node distance=0.6cm and 0.5cm,
  every node/.style={
    draw, 
    font=\tiny, 
    rounded corners=2pt, 
    inner sep=3pt, 
    align=center,
    minimum height=11pt,
  },
  arrow/.style={-Stealth}
]

\node (node1) [text width=100pt] {Requisitos - Sistema \\binario masivo};

\node (node2) [below=of node1, text width=100pt] {Origen - Supernova};

\node (node3) [below right=10pt and -20pt of node2, text width=100pt] {Colapso del núcleo};

\node (node4) [below right=30pt and -20pt of node2, text width=100pt] {Supervivencia};

\node (node5) [below=50pt of node2, text width=100pt] {Sistema binario durmiente};

\node (node6) [below=of node5, text width=100pt] {Microcuásar activo};

\node (node7) [below right=10pt and -20pt of node6, text width=100pt] {Roche-Lobe / Vientos solares};

\node (node8) [below right=30pt and -20pt of node6, text width=100pt] {Diagrama Intensidad-Dureza};

\node (node9) [below=50pt of node6, text width=100pt] {Etapas finales};

\node (node10) [below right=10pt and -20pt of node9, text width=100pt] {Fin de \textit{jets} relativistas};

\node (node11) [below right=30pt and -20pt of node9, text width=100pt] {Muerte de estrella};

\node (node12) [below right=50pt and -20pt of node9, text width=100pt] {Unifiación del sistema};

\draw [arrow] (node1.south) -- (node2.north);

\draw [arrow, color=blue!65!black] (node2.south) |- (node3.west);

\draw [arrow, color=blue!65!black] (node2.south) |- (node4.west);

\draw [arrow] (node2.south) -- (node5.north);

\draw [arrow] (node5.south) -- (node6.north);

\draw [arrow, color=blue!65!black] (node6.south) |- (node7.west);

\draw [arrow, color=blue!65!black] (node6.south) |- (node8.west);

\draw [arrow] (node6.south) -- (node9.north);

\draw [arrow, color=blue!65!black] (node9.south) |- (node10.west);

\draw [arrow, color=blue!65!black] (node9.south) |- (node11.west);

\draw [arrow, color=blue!65!black] (node9.south) |- (node12.west);

\end{tikzpicture}
\caption{Línea del tiempo representando la vida de un microcuásar.}
\end{figure}

\subsection{Requisitos para formación. Origen.}

\paragraph{Evolución de estrellas masivas y límite de Eddington.} La formación del objeto compacto, motor del microcuásar, requiere que la estrella primaria nazca en la ZAMS con una masa suficiente tal que entre en el ciclo CNO, i.e.,

\begin{equation}
  M\geq 8-10 M_{\odot}
\end{equation}

\noindent De acuerdo con las ecuaciones de homología, las estrellas masivas tienen un núcleo convectivo (debido a que $\varepsilon_{\mbox{CNO}}\propto\rho T^{12}$, que produce una luminosidad central demasiado alta para ser evacuada únicamente por radiación, violando el Criterio de Estabilidad de Schwarzschild). La envoltura, en cambio, sí es radiativa, dominada por la dispersión de Thomson en los electrones libres. De esta forma, la relación masa-luminosidad es muy elevada ($L\propto M^3$), y la estrella vive poco tiempo en la ZAMS.

\begin{equation}
  \tau_{\mbox{nuc}}\sim \frac{\varepsilon M c^2}{L}\sim 10^6-10^7\mbox{ años}.
\end{equation}

\noindent Estas estrellas, localizadas en la región de presión de radiación en el diagrama H-R, siguen la ecuación de estado

\begin{equation}
  P_{\mbox{rad}}=\tfrac{1}{3}aT^4=(1-\beta)P.
\end{equation}

\noindent Debido a su masa, el índice adiabático tenderá hacia $4/3$, el umbral de inestabilidad. Así, es necesario imponer que la fuerza de atracción sea mayor que el gradiente de la presión de radiación. Se define la luminosidad de Eddington

\begin{equation}
  L_{\mbox{Edd}}=\frac{4\pi GMc}{\kappa}\approx\frac{4\pi GMm_pc}{\sigma_T}\approx 1,3\times10^{38}\left(\tfrac{M}{M_{\odot}}\right)\approx 3,2\times 10^4\left(\tfrac{M}{M_{\odot}}\right)L_{\odot}.
\end{equation}

\noindent Esta magnitud acabará correspondiéndose con la unidad natural de luminosidad para microcuásares. 

Para que el sistema binario acabe convirtiéndose en un microcuásar, la separación orbital debe de ser suficientemente estrecha para que haya transmisión de masa desde la estrella al objeto compacto, pero sin que se fusionen al completo. 

En el caso de que el binario tenga masas muy dispares (de al menos 25 órdenes de magnitud), se deberá cumplir una fase de envoltura común: la primaria en fase de supergigante engloba a la secundaria, y la fricción hidrodinámica extrae energía orbital. La secundaria se acaba expulsando de la envolvente, y dando lugar un sistema binario compacto, con un núcleo de helio (estrella de Wolf-Rayet), y una estrella de baja masa.

Una vez agotado el hidrógeno en el centro de la primaria, se evita el \textit{flash} de helio mediante una emisión térmica de neutrinos. El núcleo encadena las siguientes reacciones:

\vspace{-0.5em}

\begin{enumerate}
  \item Combustión nuclear de $\phantom{}^4\mbox{He}$ (reacción triple-$\alpha$, $T\geq 10^8\mbox{K}$)\vspace{-0.5em}
  \item Combustión nuclear de $\phantom{}^{12}\mbox{C}$ ($T\in[5\times10^8, 10^9]K$, para una estrella de $25M_{\odot}$, $\sim600$ años)\vspace{-0.75em}
  \item Combustión nuclear de $\phantom{}^{20}\mbox{Ne}$ ($\sim 1$ año)\vspace{-0.5em}
  \item Combustión nuclear de $\phantom{}^{16}\mbox{O}$ ($T\geq1,4\times10^9\mbox{K}$, $\sim 0,5$ años)\vspace{-0.5em}
  \item Combustión nuclear de $\phantom{}^{28}\mbox{Si}$ ($T\geq 2-4\times10^{9}\mbox{K}$, $\sim 1$ día)\vspace{-0.5em}
\end{enumerate}

\noindent Al llegar a los metales férreos ($\mbox{Fe, Ni}$), con energía de ligadura muy alta ($B/A\approx8,8\mbox{MeV}$), las reacciones dejan de ser exotérmicas, y se empieza a consumir energía térmica.

Consecuentemente, el núcleo de $\phantom{}^{56}\mbox{Fe}$, originalmente sostenido por la presión de los electrones degenerados ultrarrelativistas ($\Gamma=4/3$), sigue creciendo hasta el límite de Chandrasekhar

\begin{equation}
  M_{\mbox{Ch}}=4\pi \left[\tfrac{K}{\pi G}\right]^{3/2}=5,83\mu_e^{-2}M_{\odot}\approx 1,46 M_{\odot}
\end{equation}

\noindent En este instante se produce un colapso gravitatorio. El proceso de caída libre tiene una duración aproximada de $\tau_{\mbox{libre}}\sim100\mbox{ms}$. Esta desintegración es apoyada por dos fenómenos:

\begin{enumerate}
  \item \textbf{Fotodesintegración:} $\gamma+\phantom{}^{56}\mbox{Fe}\to13\phantom{}^4\mbox{He}+4n-124,4\mbox{MeV}$, $\gamma+\phantom{}^{4}\mbox{He}+2p+2n-28,3\mbox{MeV}$. Resultado: $\Gamma\lneq 4/3$.
  \item \textbf{Captura electrónica (neutronización):} $p+e^-\to n+\nu_e$. Resultado: liberación de neutrinos, llevando potencial gravitatorio.
\end{enumerate}

\noindent Cómo finaliza este colapso indica la naturaleza del microcuásar.

\paragraph{Caso 1. \textit{Natal kick}.} Cuando la masa de las progenitoras a edad $0$ está entre $8M_{\odot}$ y $40M_{\odot}$ y $\rho\sim 2-3\rho_0$, los nucleones forman un líquido de Fermi. Éste queda gobernado por la ecuación de Tolman-Oppenheimer-Volkoff:

\begin{equation}
  \tfrac{dP}{dr}=-\frac{Gm(r)\rho(r)}{r^2}\left[1+\frac{P(r)}{\rho(r)c^2}\right]\left[1+\frac{4\pi r^3 P(r)}{m(r)c^2}\right]\left[1-2\frac{Gm(r)}{rc^2}\right]^{-1}.
\end{equation}

\noindent Las capas externas rebotan sobre el núcleo rígido, generando una onda de choque que expulsa a las exteriores, es decir, una supernova de colapso del núcleo. (En efecto, será de tipo II si conservaba $\mbox{H}$, ó Ib/Ic si no). Queda restante una estrella de neutrones $M\sim2M_{\odot}$ o un agujero negro de baja masa. Este último caso se llama \textit{fallback}, y ocurre cuando la caída del material supera el límite de TOV.

Apelando a la <<Arqueología Galáctica>> \cite{mirabel}, la eyección de masa mediante el efecto Blaauw y acumulación de asimetrías hidrodinámicas conllevan un fuerte impulso al centro de masas, conocido como \textit{Natal kick}. 


\paragraph{Caso 2. Colapso directo.} Cuando la masa en la ZAMS supera $25-40M_{\odot}$ y se tiene una baja metalicidad, el núcleo de hierro supera directamente el límite de estabilidad de TOV. La onda de choque fracasa y la estrella implosiona sin producir una supernova luminosa ni expulsar una cantidad de masa significativa.

Los sistemas binarios resultantes tienen velocidades extremadamante bajas, normalmente asociados a rotación galáctica. Sus agujeros negros fueron formados por colapso directo. 

Si la estrella tuviera momento angular previo, se formarán tanto un agujero negro como un disco de acreción ultradenso, expulsando chorros ultrarrelativistas que perforan la estrella primaria. Ésta es perforada, y se produce un estallido de rayos gamma de larga duración.


\subsection{Sistema binario durmiente}

\noindent Si el sistema binario se ha formado de manera correcta y no colapsa, se puede permanecer una cantidad de tiempo indeterminada (y posiblemente muy larga) en un estado durmiente (también conocido como estado de quiesencia). En el caso particular de las binarias de rayos X de rayos X (LMXB), esta será la etapa más larga de su vida útil. 

Los episodios de alta intensidad (\textit{outbursts}) duran unos meses, y van intercalados por intervalos de recurrencia de años o décadas. En este periodo, el ritmo de transferencia es extremadamente bajo, 

\begin{equation}
  L_X\sim 10^{-8,5}-10^{-6,5}L_{\mbox{Edd}}.
\end{equation}

\noindent Según Malzac \cite{malzac}, el disco de acreción se halla a una gran distancia del agujero negro, mientras que la región interna contiene un flujo caliente advectivo tenue e ineficiente. (ADAF, RIAF).

Para activar un microcuásar durmiente, uno debe referirse al modelo de Meyer-Hofmeister, basado en la inestabilidad térmica-viscosa asociada a la ionización del hidrógeno en el disco de acreción. Los pasos a seguir son:

\begin{enumerate}
  \item El gas transferido desde la secundaria hasta la primaria se acumula en las regiones externas del disco. El hidrógeno está en el estado neutro de la ecuación de Saha. La viscosidad efectiva es minúscula, el flujo entrante sobrepasa al cebado del núcleo.
  \item Según la masa se acumula en el exterior, crece la densidad y temperatura. Entorno a los $6000$ ó $10000$ Kelvin, se ioniza el hidrógeno.
   
  En este régimen, la opacidad de Rosseland aumenta (debido a $\mbox{H}^-$ y las transiciones \textit{bb}). Además, el índice adiabático ($\Gamma=\partial_{\ln\rho}\ln P$) cae: $4/3\to 1\mbox{ o }2$, provocando inestabilidad convectiva.
  \item Cuando $T\geq 10^4\mbox{K}$, el gas está totalmente ionizado. El acoplamiento magnético del cono de plasma dispara la viscosidad. Se desencadena una avalancha de materia, aumentando la tasa de acreción y la luminosidad. 
\end{enumerate}


\begin{figure}[h]
\centering
\makebox[\textwidth][c]{
\begin{tikzpicture}[>=Stealth, font=\sffamily]

\begin{scope}[xshift=0cm]
    \node[above=1.4cm, font=\bfseries] at (-0.3,0) {Binaria primordial};
    
    \draw[dash pattern=on 2.5pt off 2pt, line width=0.6pt] 
      (-0.45, 0) to[out=135, in=0] (-1.2, 0.6)
      arc (90:270:0.6)
      to[out=0, in=225] (-0.45, 0);

    \draw[dash pattern=on 2.5pt off 2pt, line width=0.6pt] 
      (-0.45, 0) to[out=45, in=180] (0.5, 1.2)
      arc (90:-90:1.2)
      to[out=180, in=-45] (-0.45, 0);

    \fill (-1.2, 0) circle (2.2pt);
    \fill (0.5, 0) circle (6.5pt);
  \end{scope}

  \draw[->, line width=0.7pt] (1.9, 0) -- (2.6, 0);

  \begin{scope}[xshift=4.6cm]
    \node[above=1.4cm, font=\bfseries] at (-0.3,0) {Roche-lobe};
    
    \fill[shading=diffusestar] (0.5, 0) circle (1.1);

    \draw[dash pattern=on 2.5pt off 2pt, line width=0.6pt] 
      (-0.45, 0) to[out=135, in=0] (-1.2, 0.6)
      arc (90:270:0.6)
      to[out=0, in=225] (-0.45, 0);

    \draw[dash pattern=on 2.5pt off 2pt, line width=0.6pt] 
      (-0.45, 0) to[out=45, in=180] (0.5, 1.2)
      arc (90:-90:1.2)
      to[out=180, in=-45] (-0.45, 0);

    \fill (-1.2, 0) ellipse (0.38 and 0.08);
    \fill (-1.2, 0) circle (1.8pt);
    \draw[line width=0.7pt] (-0.45, 0) to[out=180, in=0] (-0.85, -0.02);
    \fill (0.5, 0) circle (6.5pt);
  \end{scope}

  \draw[->, line width=0.7pt] (6.7, 0) -- (7.4, 0);

  \begin{scope}[xshift=9.2cm]
    \node[above=1.4cm, font=\bfseries] at (0,0) {Espiral};
    
    \fill[shading=diffusestar] (0, 0) circle (1.15);
    \fill (0, 0) circle (7pt);
    \draw[line width=0.6pt] (0, 0) ellipse (0.85 and 0.8);
    \fill (-0.5, 0.65) circle (2.6pt);
  \end{scope}

  \draw[->, line width=0.7pt] (10.65, 0.35) -- (11.35, 0.85);
  \draw[->, line width=0.7pt] (10.65, -0.35) -- (11.35, -0.85);

  \begin{scope}[xshift=11.9cm, yshift=0.9cm]
    \begin{scope}[scale=0.6]
      \draw[line width=0.6pt] 
        (-0.25, 0) to[out=135, in=0] (-0.7, 0.35)
        arc (90:270:0.35)
        to[out=0, in=225] (-0.25, 0);

      \draw[line width=0.6pt] 
        (-0.25, 0) to[out=45, in=180] (0.4, 0.65)
        arc (90:-90:0.65)
        to[out=180, in=-45] (-0.25, 0);

      \fill (-0.7, 0) circle (2.2pt);
      \fill (0.4, 0) circle (4.5pt);
    \end{scope}
    \node[anchor=west, font=\small\bfseries] at (0.7, 0) {Binario compacto};
  \end{scope}

  \begin{scope}[xshift=11.9cm, yshift=-0.9cm]
    \fill (0.1, 0) circle (3pt);
    \node[anchor=west, font=\small\bfseries] at (0.7, 0) {Estrella aislada};
  \end{scope}

\end{tikzpicture}}
\caption{Recreado desde XiangDong 2010.}
\end{figure}

\vspace{1em}

\subsection{Etapa activa}

\noindent Se refiere al lector a la sección siguiente.



\subsection{Desmantelamiento sistémico}

\noindent La etapa activa concluye cuando el suministro de materia de la estrella secundaria acaba. La estructura disco-corona-\textit{jet} se desmantela sistémicamente. Existen dos mecanismos posibles para esto, dependiendo de la masa restante.

\begin{enumerate}
  \item \textbf{LMXB:} La estrella donante pierde su envoltura de $\mbox{H}$. El núcleo remanante se enfría como una enana blanca, dejando un sistema binario pasivo: una enana blanca y un agujero negro/púlsar.
  \item \textbf{HMXB:} En binarias de alta masa (p.ej. Cygnus X-1, Cygnus X-3), la estrella donante (O/B/Wolf-Rayet) agota su combustible nulcear y se colapsa el núcleo por la fuerza gravitatoria. Si la pérdida de masa o el segundo \textit{Natal Kick} (esta vez asimétrico) supera la energía de ligadura, el sistema presenta una disrupción. Ambos compactos son disparados al espacio interestelar, en distintos caminos con altas velocidades. Si no se desligan, se forma un binario de dos objetos compactos: NS-NS (PSR 1913+16), púlsar doble (J0737-3039), NS+BH ó BH+BH.
\end{enumerate}

\noindent A partir de este momento, la evolución final es gobernada por la relatividad general. El momento cuadripolar de ambas masas compactas,

\begin{equation}
  Q_{ij}=\int\rho(x)(x_ix_j-\tfrac{1}{3}r^2\delta_{ij})dx
\end{equation}

\noindent da lugar a ondas gravitacionales. Su potencia cuadripolar es

\begin{equation}
  L_{GW}=\tfrac{1}{5}\tfrac{G}{c^5}\left\langle \dddot{Q}_{ij}, \dddot{Q}^{ij}\right\rangle\sim \frac{c^5s^2}{G}\left(\frac{R_S}{R}\right)^2\left(\frac{v}{c}\right)^6.
\end{equation}

\noindent La pérdida de energía conlleva una reducción en el radio orbitial. El periodo orbital estará gobernado por

\begin{equation}
  \Delta T(t)=\frac{\dot{P}_{\mbox{orb}}}{2P_{\mbox{orb}}}t^2
\end{equation}

\noindent en un tiempo $\tau\sim10^7-10^9$ años. 

Las fases finales (\textit{inspiral}, \textit{merger}, \textit{ringdown}) son detectadas en la Tierra por interferómetros como LIGO.


% Véase ACCRETION TURNOFF AND RAPID EVAPORATION OF VERY LIGHT SECONDARIES IN LOW-MASS X-RAY BINARIES M. Ruderman, J. Shaham, and M. Tavani
